\input{estilo}
% Propuesta: el polling en el tiempo. El banco pregunta cada Δt; la mitad de
% las respuestas dicen que nada cambió, y un precio que duró menos que Δt no lo
% ve nadie. Agregado: los valores, el 103 perdido y el "× símbolos × bancos",
% que salen del texto y no de la pizarra.
\begin{document}
\begin{figura}\begin{tikzpicture}
  \def\yn{2.2}
  % las dos líneas de tiempo
  \foreach \y/\n in {\yn/NASDAQ, 0/banco} {
    \node[anchor=east] at (-0.25,\y) {\n};
    \draw[->, draw=suave, line width=0.5pt] (0,\y) -- (11.2,\y);
  }
  \node[etiq, text=suave, anchor=south west] at (0.1,\yn) {\var{GOOGL} = 100};
  % los cambios de precio en el NASDAQ
  \foreach \x/\v in {4.7/103, 5.6/101} {
    \node[evento=acento] at (\x,\yn) {};
    \node[etiq, text=acento, above=2pt] at (\x,\yn) {\v};
  }
  \node[etiq, text=acento, anchor=south] at (4.7,\yn+0.45) {no lo ve nadie};
  % las consultas, cada Δt, y lo que vuelve
  \foreach \x/\v/\c [count=\i] in {1.0/100/, 3.6/100/sin cambios, 6.2/101/, 8.8/101/sin cambios} {
    \draw[msg] (\x,0) -- (\x+0.4,\yn);
    \draw[respuesta] (\x+0.55,\yn) -- (\x+0.95,0);
    \node[evento] at (\x,0) {};
    \node[etiq, below=3pt] at (\x+0.95,0) {\v};
    \node[etiq, text=suave, below=16pt] at (\x+0.95,0) {\c};
  }
  % el período
  \draw[llave=5pt] (3.6,-1.25) -- (1.0,-1.25);
  \begin{scope}[shift={(2.05,-1.85)}]
    \draw (0,0) circle (2.2mm);
    \draw (0,0) -- (0,0.15) (0,0) -- (0.1,-0.04);
  \end{scope}
  \node[etiq, anchor=west] at (2.3,-1.85) {Δ\var{t}};
  \node[etiq, text=suave, anchor=east] at (11.2,-1.85) {y todo esto, por cada símbolo y por cada banco};
\end{tikzpicture}\end{figura}
\end{document}
